and define a homomorphism:
@@M0@@
by putting:
@@M1@@
with @@M2@@.

\sgaStableTarget{sga6.e07.prop.4-7-5}\textbf{Proposition 4.7.5.} \uline{The operations defined by} (4.7.1) \uline{and} (4.7.2) \uline{endow} @@M3@@ \uline{with a} @@M4@@-\uline{ring structure}, \uline{and} @@M5@@ \uline{is an injective homomorphism identifying} @@M6@@, \uline{equipped with the operations defined by} (4.7.3) \uline{and} (4.7.4), \uline{with a} @@M7@@-\uline{ideal of} @@M8@@.

It is immediate to verify that @@M9@@ is a @@M10@@-ring: one can refer, for example, to V 5.5, where an entirely analogous construction is given. The second part of the proposition is also proved without difficulty using the preceding relations, the relation @@M11@@ which follows immediately from the hypotheses, the relation @@M12@@ (see VI 5.9), and the relations of V 5.

Put then:
@@M13@@
and denote by @@M14@@ the canonical homomorphism @@M15@@. Regard @@M16@@ as a @@M17@@-ring through the quotient structure.

Next define additive homomorphisms:
@@M18@@

The homomorphism @@M19@@ is the composite @@M20@@. Likewise, the homomorphism @@M21@@ is the composite @@M22@@. We then have:
@@M23@@
